\printedpage{194}
% head: A. Grothendieck : Dérivateurs, Ch. XVIII, Édition des Universités de Montpellier II et Paris VII

\noindent\hspace*{1em}\textcircled{a}\quad $\chi(f_A)+\chi(f_B)=1$, and\\
\hspace*{1em}\textcircled{b}\quad $\chi(f_A)\chi(f_B)=0$.\\
$\chi(f_A)\neq 1$, we conclude from \textcircled{a} $\chi(f_B)\neq 0$, whence by \textcircled{b} $\chi(f_A)=0$, whence from \textcircled{a} $\chi(f_B)=1$, whence $B\in\Phi$.

\msp{page 7}
Let us finally prove that $\mathfrak{m}=\mathfrak{m}_\Phi$. Since $\mathfrak{m}_\Phi$ is a maximal ideal, it suffices to prove
\[
\mathfrak{m}_\Phi\ \subset\ \mathfrak{m}\ .
\]
Now if $A\in\mathfrak{P}(E)$ is such that $A\notin\Phi$, or $f_A\in\mathfrak{m}_\Phi$ (cor.~1), we also have $f_A\in\mathfrak{m}$. Thus it suffices to prove the

\textbf{Lemma.} \textit{Let $\Phi\in\mathcal{F}_0$. Then $\mathfrak{m}_\Phi$ is the ideal generated by the $f_A$ for $A\in\mathfrak{P}(E)\backslash\Phi$. More precisely, the $f_A$ are in $\mathfrak{m}_\Phi$, and every $f\in\mathfrak{m}_\Phi$ is of the form $ff_A$, with $A$ as above.}

Indeed, $f\in\mathfrak{m}_\Phi$ means that there exists $B\in\Phi$ such that $f|B=0$. But this also means that $f=ff_A$, where $A=E\backslash A$ \texttt{[rather} $A=E\backslash B$\texttt{]}, whence $A\notin\Phi$.

This proves proposition 2, and at the same time:

\textbf{Corollary 2.} \textit{Every quotient ring of $k_E=\prod_{\xi\in E}k_\xi$ which is an integral domain is a field.}

Let us now note that if $\Phi\in\mathcal{F}$ is any filter on $E$, then $k_\Phi$ is a quotient ring of $k_E$, hence
\[
\operatorname{Max}(k_\Phi)\ \hookrightarrow\ \underbrace{\operatorname{Max}(k_E)}_{\simeq\,\mathcal{F}_0}\ . \tag*{$(*)$}
\]
This being so, we find:

% a vertical bar is printed in the left margin beside the statement of Corollaire 3
\textbf{Corollary 3} (\footnote[206]{Works for any system of rings $(k_\xi)_{\xi\in E}$ \textit{which are nonzero}, and $\Phi,\Psi\in\mathcal{F}$, $\Psi$ finer than $\Phi$, i.e. $\Phi\subset\Psi$ if and only if $\mathfrak{m}_\Phi\subset\mathfrak{m}_\Psi$ (no need for $\Psi$ to be an ultrafilter).}). \textit{The preceding inclusion identifies} $\operatorname{Max}(k_\Phi)$ \textit{with the set of ultrafilters $\Psi$ on $E$ which are} finer \textit{than $\Phi$.}

Indeed, if $\Psi$ is such an ultrafilter, we
\msp{page 8}
have the commutative diagram
\[
\begin{tikzcd}
k_E \arrow[r, "\chi_\Phi"] \arrow[rr, bend right=15, "\chi_\Psi"'] & k_\Phi \arrow[r] & k_\Psi
\end{tikzcd}\ , \tag*{$(**)$}
\]
which shows that $\Psi$ \texttt{[rather} $\mathfrak{m}_\Psi$\texttt{]} is in the image of $(*)$. Conversely, suppose that $\Psi\in\mathcal{F}_0$ comes from $\operatorname{Max}(k_\Phi)$, i.e. that we have $\mathfrak{m}_\Phi\subset\mathfrak{m}_\Psi$ (where $\mathfrak{m}_\Phi$ is the kernel of $\chi_\Phi$), or again that we have a factorization $(**)$ of $\chi_\Psi$. Let us prove that $\Psi$ is finer than $\Phi$, i.e. that $\Psi\supset\Phi$. So let $A\in\Phi$, let us prove that we have $A\in\Psi$. But $A\in\Phi$ implies \texttt{[the]} fact that $\chi_\Phi(f_A)=1$, hence its image in $k_\Psi$ is equal to 1, hence we have $A\in\Psi$ by corollary 1.

\printedpage{195}
% head: A. Grothendieck : Dérivateurs, Ch. XVIII, Édition des Universités de Montpellier II et Paris VII

\msp{page 9}
Let us now suppose that all the $k_\xi$ are equal \texttt{[to]} one and the same ring \fbox{$\ k\neq 0\ $}\,. Thus $k_J$, $k_\Phi$ are ``rings under $k$'' $k\longrightarrow k_J\longrightarrow k_\Phi$, moreover $k\longrightarrow k_J$ (if $J\neq\emptyset$) and $k\longrightarrow k_\Phi$ are injective. We ask the question when $k\longrightarrow k_\Phi$ is \textit{bijective}, i.e. surjective. Since $k_\Phi\simeq k_E/\mathfrak{m}_\Phi$, this therefore means that $k_E=k+\mathfrak{m}_\Phi$, i.e. that every $f\in k_E$ can be written
\[
f\ =\ c\ +\ g\qquad
\begin{cases}
c\ \in\ k\ , & \text{i.e. } c(\xi)=\underbrace{c_0}_{\in\,k}\ \forall\,\xi\in E\ ,\\[2ex]
g\ \in\ \mathfrak{m}_\Phi & \text{i.e. } \exists\,A\in\Phi \text{ with } g|A=0\ .
\end{cases}
\]
But for $f\in k_E$, $c\in k$ and $A\in\Phi$ fixed, the existence of $g\in k_E$ such that $g|A=0$ and $f=c+g$ is equivalent to $f|A=c_A$ (constant function on $A$, with value $c$), i.e. $A\subset f^{-1}(\{c\})$. For $f$, $c$ fixed, the existence of $A\in\Phi$ satisfying the relation is equivalent to $f^{-1}(\{c\})\in\Phi$. And for $f$ fixed, the existence of $c$ is equivalent to that of a fiber $E_c$ of $f:E\longrightarrow k$, which is in $\Phi$. Thus:

\textbf{Proposition 3.} \textit{Let $\Phi\in\mathcal{F}$. For $k\longrightarrow k_\Phi$ to be surjective, it is necessary and sufficient that for every family of subsets} essentially disjoint \textit{$(E_\xi)_{\xi\in k}$ of $E$ indexed by $k$, with union $E$, there exist a $\xi\in k$ such that $E_\xi\in\Phi$. A fortiori, it} is necessary \textit{that $\Phi$ be an ultrafilter (since} $\operatorname{card}k\geqslant 2$ \textit{by hypothesis).}

\msp{page 10}
Indeed, giving such a family $(E_\xi)_{\xi\in k}$ is equivalent to giving a map $E\longrightarrow k$, i.e. an element of $f_E$. % [sic] $f_E$ -> $k_E$, an application $E\to k$ is an element of $k_E$

\textbf{Scholie.} The property we have just made explicit for a filter $\Phi$ on $E$, depends on $k$ only through $\pi=\operatorname{card}k$. We can also state it by saying that for every \textit{partition} $(E_i)_{i\in I}$ of $E$ into a set $I$ of subsets such that $\operatorname{card}I\leqslant\pi$, there is one (and only one) of the $E_i$ which is in $\Phi$. We then say that $\Phi$ is a $\pi$-\textit{ultrafilter.}

If $\pi\geqslant 2$ is finite, this property is equivalent to saying that $\Phi$ is an ultrafilter. Besides, whatever $\pi$, the ultrafilters $\Phi$ \textit{which are trivial} (i.e. such that there exists $\xi\in E$ with $\{\xi\}\in\Phi$, or again such that $\Phi=\{A\in\mathfrak{P}(E)\mid\xi\in A\}$) are $\pi$-ultrafilters. If $\pi$ is infinite, it is not clear that there exist nontrivial $\pi$-ultrafilters. We know that such exist on $E$ \textit{only} if $\operatorname{card}E\notin\mathfrak{U}_1$, where $\mathfrak{U}_1$ is the smallest universe such that $\mathbf{N}\in\mathfrak{U}_1$, i.e. such
\msp{page 11}
that $\mathfrak{U}$ contains as an element an infinite set.

If $\mathfrak{U}=\mathfrak{U}_1$, and if $k$ is an infinite field, e.g. $k=\mathbf{Q}$, then for every set $E\in\mathfrak{U}$ and for every \textit{nontrivial} ultrafilter $\Phi$ on $E$, $k_E$ is a \textit{nontrivial} extension of $k$. But if $\mathfrak{U}=\mathfrak{U}_\alpha$ is an arbitrary universe, and $\pi$ the first cardinal in $\mathfrak{U}$ which does not belong to a strictly smaller universe $\mathfrak{U}'\subset\mathfrak{U}$, I presume that for every $E\in\mathfrak{U}$, every $\pi$-ultrafilter on $E$ is trivial. % [sic] $k_E$ est une extension -> $k_\Phi$ est une extension, the extension of $k$ attached to the ultrafilter $\Phi$ is $k_\Phi$

\printedpage{196}
% head: A. Grothendieck : Dérivateurs, Ch. XVIII, Édition des Universités de Montpellier II et Paris VII

\subsection*{Application to an example of a nonaccessible subcategory.}

Let $\mathcal{M}$ be a $\mathfrak{U}$-accessible category, and $a\in\operatorname{Ob}\mathcal{M}$, whence a functor
\[
\begin{aligned}
f^o\ :\ \mathcal{M}^o\ &\longrightarrow\ \mathrm{Ens}\\
x\ &\longmapsto\ \operatorname{Hom}_{\mathcal{M}}(x,a)\ \overset{\text{def}}{=}\ a(x)\ ,
\end{aligned}
\]
a functor commuting with arbitrary $\limproj$, hence such that $f:\mathcal{M}\longrightarrow(\mathrm{Ens})^o$ commutes with arbitrary $\limind$, and in particular $f$ is accessible. Let us consider the strictly full subcategory
\msp{page 12}
\[
\underbrace{\mathrm{Ens}^*}_{=\,\mathrm{Ens}\backslash\{\emptyset\}}\ \subset\ \mathrm{Ens}
\]
of nonempty sets, we know that it is accessible, but it is not stable under $\limproj$ of type $\mathrm{L}_\pi$, however large we take $\pi$. % [sic] n'est par stable -> n'est pas stable, typo

Let us consider
\[
\underbrace{\underline{S}\ =\ \underline{S}_a}_{\overset{\text{def}}{=}\,f^{-1}(\mathrm{Ens}^*)}\ \subset\ \mathcal{M}\ ,
\]
strictly full subcategory of $\mathcal{M}$ consisting of the $x\in\operatorname{Ob}\mathcal{M}$ such that
\[
a(x)\ \neq\ \emptyset\ ,\qquad\quad \text{i.e. } x\leqslant a \text{ for the canonical preorder of } \mathcal{M}\ .
\]
It is therefore an \textit{open} subcategory of $\mathcal{M}$. Its closed complement $\underline{S}'=\underline{S}'_a$ is the closed subcategory which is the inverse image, under $f$, of the strictly full subcategory of Ens reduced to the single object $\{\emptyset\}$. This subcategory of Ens is the standard final category. It is stable in Ens under arbitrary $\limind$, and under arbitrary filtering $\limproj$. Thus, $\underline{S}'$ is stable in $\mathcal{M}$ under arbitrary filtering $\limind$. Thus, $\underline{S}'$ is stable in $\mathcal{M}$ under arbitrary filtering $\limind$, in particular, it is $\mathrm{L}_\pi$-stable for $\pi$ sufficiently large. Let us make this explicit:

\msp{page 13}
\[
\begin{array}{cl}
\underbrace{\underline{S}'\ =\ \underline{S}'_a}_{}\ \subset\ \mathcal{M} & \\
\Big\|\ \text{\scriptsize def} & \text{closed subcategory of } \mathcal{M}\\
f^{-1}(\{\emptyset\})\ = & \text{consisting of the } x\ \in\ \operatorname{Ob}\mathcal{M} \text{ such}\\
 & \text{that } a(x)=\emptyset.
\end{array}
\]

Since $\underline{S}'$, being closed in $\mathcal{M}$, is stable under nonempty $\limind$ and in particular $\mathrm{L}_\pi$-stable for $\pi$ large, there is more reason to expect $\underline{S}'$ to be accessible in $\mathcal{M}$ than for its complement $\underline{S}$, for which it is not clear that it has this stability property, but only stability under nonempty $\limproj$, and in particular under direct factors.

In fact, we shall see that it can happen that $\underline{S}$ and $\underline{S}'$ are both nonaccessible in $\mathcal{M}$:

\textbf{Proposition 4.} \textit{Let $\mathfrak{U}$ be a universe, containing an infinite cardinal $\pi_0$ such that for every $E\in\mathfrak{U}$, every $\pi_0$-ultrafilter on $E$ is trivial. Let then $k\in\mathfrak{U}$ be a field of cardinality $\pi'_0\geqslant\pi_0$, and let $\mathcal{M}$ be the category}
\printedpage{197}
% head: A. Grothendieck : Dérivateurs, Ch. XVIII, Édition des Universités de Montpellier II et Paris VII
\msp{page 14}
\textit{of commutative $k$-algebras. Thus $\mathcal{M}$ is a $\mathfrak{U}$-paratopos, admitting $k$ as initial object. Let $\underline{S}'$ be the strictly full subcategory, stable under small nonempty $\limind$, consisting of the objects $A$ of $\mathcal{M}$ such that $\operatorname{Hom}_{\mathcal{M}}(A,k)=\emptyset$, and $\underline{S}$ the complementary strictly full subcategory, consisting of the objects $A$ such that $\operatorname{Hom}_{\mathcal{M}}(A,k)=\emptyset$. Then $\underline{S}'$ is a subcategory} closed \textit{in $\mathcal{M}$, and $\underline{S}$ the complementary open subcategory, hence $\underline{S}$ (just as $\underline{S}'$) is stable under direct factors. This being so:} % [sic] second $\operatorname{Hom}_{\mathcal{M}}(A,k)=\emptyset$ -> $\operatorname{Hom}_{\mathcal{M}}(A,k)\neq\emptyset$, $\underline{S}$ is the complement of $\underline{S}'$
\begin{enumerate}[label={}]
\item[\textit{a)}] \textit{There is no cardinal $\pi$ such that $\underline{S}$ is $L_\pi$-stable in $\mathcal{M}$, a fortiori $\underline{S}$ is not accessible in $\mathcal{M}$.}
\item[\textit{b)}] (\footnote[207]{?, not proved.}). \textit{Although $\underline{S}'$ is $L_0$-stable in $\mathcal{M}$ (i.e. stable under filtering $\limind$, and even, as we have said, under every nonempty $\limind$), $\underline{S}'$ is nevertheless not accessible in $\mathcal{M}$.}
\end{enumerate}
\noindent\textsc{Proof}. Assertion a) will follow from the

\textbf{Corollary 1.} \textit{Under the conditions of proposition 4, let $E\in\mathfrak{U}$ be an infinite set, $\Phi_0$ the filter of complements of finite subsets of $E$. Then $\operatorname{Hom}_{\mathcal{M}}(k_{\Phi_0},k)=\emptyset$, a fortiori, for every filter $\Phi$ on $E$ finer than $\Phi_0$, we have $\operatorname{Hom}_{\mathcal{M}}(k_\Phi,k)=\emptyset$.}

\msp{page 15}
Here we have put, of course, for every filter $\Phi$ on $E$,
\[
k_\Phi\ =\ \limind_{J\in\Phi}k_J\ .
\]
We know that for every commutative $k$-algebra $\mathcal{A}$, $\operatorname{Hom}(\mathcal{A},k)$ is in 1-1 correspondence with the subset of $\operatorname{Max}(\mathcal{A})$ consisting of the maximal ideals $\mathfrak{m}$ of $\mathcal{A}$ such that $k\longrightarrow\mathcal{A}/\mathfrak{m}$ is an isomorphism. If $\mathcal{A}$ is of the form $k_\Phi$, $\Phi$ a filter on $E$, this set is in one-to-one correspondence with the set of $\pi'_0$-ultrafilters $\Psi$ on $E$ finer than $\Phi$, where $\pi'_0=\operatorname{card}(k)$ (proposition 3, and cor.~3 of prop.~2, page 7). Since $\pi'_0\geqslant\pi_0$, such an ultrafilter is trivial by hypothesis. Thus we find
\[
\operatorname{Hom}_{\mathcal{M}}(k_\Phi,k)\ \simeq\ \text{set of } \xi\in E \text{ such that } \{\xi\}=\textstyle\bigcap_{A\in\Phi}A\ . \tag*{$(*)$}
\]
But saying that $\Phi$ is finer than $\Phi_0$ also means that for every $\xi\in E$, we have $E\backslash\{\xi\}\in\Phi$, i.e. $\exists\,A\in\Phi$ with $A\subset E\backslash\{\xi\}$, i.e. $\{\xi\}\notin\bigcap_{A\in\Phi}A$. Hence we then have $\operatorname{Hom}_{\mathcal{M}}(k_\Phi,k)=\emptyset$.

The conclusion of corollary 1 is expressed by the relation
\[
k_\Phi\ \notin\ \operatorname{Ob}\underline{S}\ ,\qquad\quad \text{i.e. } k_\Phi\in\operatorname{Ob}\underline{S}'\ ,
\]
on the other hand, we have by definition
\[
k_\Phi\ =\ \limind_{J\in\Phi}k_J\ ,
\]
\msp{page 16}
and it is clear (since these $J\in\Phi$ are $\neq\emptyset$), that the $k_J$ in question are in $\underline{S}$. Thus, for a filter $\Phi$ as in corollary 1, we find a direct system of objects of $\underline{S}$, indexed by $\Phi^o$, whose direct limit is not in $\underline{S}$. Thus a) of proposition 4 follows by noting that for every cardinal $\pi\in\mathfrak{U}$, we can find a set $E\in\mathfrak{U}$, and a filter $\Phi$ on $E$ finer than $\Phi_0$, such that the ordered set $\Phi^o$ is large relative to $\pi$. To find such a $\Phi$, it suffices to take $E$ of cardinality $\pi'>\pi$, and for $\Phi$ the set
\printedpage{198}
% head: A. Grothendieck : Dérivateurs, Ch. XVIII, Édition des Universités de Montpellier II et Paris VII
of complements of subsets of $E$ of cardinality $\leqslant\pi$, so that $\Phi^o$ is isomorphic to the ordered set of subsets of $E$ of cardinality $\leqslant\pi$, which is indeed large relative to $\pi$. This proves a) of proposition 4.

Let us now prove b). Let us note the

\textbf{Lemma.} \textit{Let $\mathcal{M}$ be an accessible category stable under $\limind$ of type $I_\alpha$, $\alpha$ a given limit ordinal, $\underline{S}'$ a closed subcategory, $\pi$ a useful cardinal}
\msp{page 17}
\textit{for $\mathcal{M}$, such that} $\operatorname{card}\alpha\leqslant\pi$\textit{. For every $x\in\operatorname{Ob}\underline{S}'$, let us choose an $\omega$-$\pi$-adapted representation ($\omega$ being the first infinite ordinal)}
\[
x\ \simeq\ \limind_{i\in I_x}x_i
\]
\textit{(thus $I_x$ large relative to $\pi$, stable under $\limind$ of type $I_\alpha$, $i\longmapsto x_i:I\longrightarrow\mathcal{M}$ commutes with the said limits, finally the $x_i\in\operatorname{Ob}\mathcal{M}_\pi$). For $\underline{S}'$ to be $\pi$-accessible in $\mathcal{M}$, it is necessary and sufficient that for every $x\in\operatorname{Ob}\underline{S}'$, there exist an $i\in I_x$ such that $x_i\in\operatorname{Ob}\underline{S}'$.}

Indeed, note that $\underline{S}'$, being closed, is stable in $\mathcal{M}$ under every nonempty $\limind$ representable in $\mathcal{M}$, and in particular under $\limind$ of type $I_\alpha$. So we can apply 4.13.22 (p.~342) to $\mathcal{M}$, $\underline{S}'$, $\pi$ and to $\Phi=\{I_\omega\}$. Then a) $\Longrightarrow$ c) implies that if $\underline{S}'$ is $\pi$-accessible in $\mathcal{M}$, then for every $x\in\operatorname{Ob}\underline{S}'$, the set $I_{x,S'}$ of the $i\in I_x$ such that $x_i\in\operatorname{Ob}\underline{S}$ is cofinal in $I_x$. This implies that it is nonempty. (And conversely too, $\underline{S}'$ being a closed subcategory in $\mathcal{M}$, since if $x_i\in\underline{S}'$, then for $j\in\operatorname{Ob}I_x$ with $i\leqslant j$, we have $x_j\in\underline{S}'$, i.e. $I_{x,\underline{S}'}$ contains the essential image of $i\backslash I_x$, which is cofinal in $I$ \dots) This proves the ``only if''. Let us prove the ``if''. So let us prove that $\underline{S}'$ is equal to the essential image of $\operatorname{Ind}_\pi(\underline{S}'(\pi))$. It contains it, since $\underline{S}'$ is stable under $\limind$ of type $\mathrm{L}_\pi$. % Doubt resolved against PDF leaf 40: the first S-prime subscript is not underlined, retained. % [sic] $x_i\in\operatorname{Ob}\underline{S}$ -> $x_i\in\operatorname{Ob}\underline{S}'$, the set $I_{x,\underline{S}'}$ concerns the $x_i$ in $\underline{S}'$

\msp{page 18}
So it remains to prove that every $x\in\operatorname{Ob}\underline{S}'$ is in the essential image. Now the hypothesis assures us that $x\simeq\limind\limits_{i\in I_x}x_i$ is isomorphic to $\limind\limits_{i\in I_{x,\underline{S}'}}x_i$ (for $I_{x,\underline{S}'}$ is cofinal in $I_x$), where this time the $x_i$ are in $\underline{S}'$. As $I_x$ large relative to $\pi$, likewise the cofinal full subcategory $I_{x,\underline{S}'}$, so we win.

\textbf{Corollary.} \textit{Under the preceding conditions, if $\underline{S}'$ is $\pi$-accessible, the same is true of the complementary open subcategory $\underline{S}$,} provided \textit{that $\underline{S}$ is stable in $\mathcal{M}$ under $\limind$ of type $L_\pi$.}

We must prove that $\underline{S}$ is equal to the essential image of $\operatorname{Ind}_\pi\underline{S}(\pi)$. The ``provided'' hypothesis ensures that this essential image is $\subset\underline{S}$. Conversely, every [\,?] $x\in\operatorname{Ob}\underline{S}$ is in this essential image. We will have
\[
x\ \simeq\ \limind_{i\in I}x_i\ ,\qquad\quad x_i\in\operatorname{Ob}\mathcal{M}_\pi,\ I=\mathcal{M}_\pi/x,\ \text{large relative to } \pi\ ,
\]
and since $x\in\operatorname{Ob}\underline{S}$ and $\underline{S}$ open, the $x_i$ are in $\underline{S}$, hence in $\underline{S}(\pi)$, OK.

In the case which concerns us, where $\mathcal{M}=(k\text{-alg.})$, as $\underline{S}$ is precisely \textit{not} stable in $\mathcal{M}$ under $\mathrm{L}_\pi$ (whatever $\pi$) (that is what made us conclude that
\msp{page 19}
$\underline{S}$ is not accessible in $\mathcal{M}$), we therefore cannot infer from this inaccessibility of $\underline{S}$, that of $\underline{S}'$. Let us note the